% answers_short.tex —— 2020 年·全国Ⅰ卷·文科 · 简短答案速查
% 由 tools/make_short_answers.py 自答案要点 MD 逐题答案行生成，勿手改；
% 选择题字母经三源交叉核对（逐题行 / 速查表 / 解析 \gkanswer 框）。

\begin{gknotice}
\textbf{使用说明}：本节只给最终答案便于快速核对，完整推导见解析版。答案为编写者独立推导，\textbf{非官方标准答案}。
\end{gknotice}

\gkpart{选择题}
\begin{center}\begin{tabular}{ccccccccccccc}
\toprule
题号 & 1 & 2 & 3 & 4 & 5 & 6 & 7 & 8 & 9 & 10 & 11 & 12 \\
\midrule
答案 & D & C & C & A & D & B & C & B & C & D & B & A \\
\bottomrule
\end{tabular}\end{center}

\gkpart{填空题}
\begin{itemize}[leftmargin=2.6em]
  \item[\textbf{13.}] $1$
  \item[\textbf{14.}] $5$
  \item[\textbf{15.}] $y=2x$
  \item[\textbf{16.}] $7$
\end{itemize}

\gkpart{解答题}
\begin{itemize}[leftmargin=2.6em]
  \item[\textbf{17.}] (1) 甲 $0.4$，乙 $0.28$；(2) 甲平均利润 $15$ 元、乙 $10$ 元，应选甲分厂
  \item[\textbf{18.}] (1) 面积 $=\sqrt3$；(2) $C=15^\circ$
  \item[\textbf{19.}] (1) 见解析版；(2) $V=\dfrac{\sqrt6}8$
  \item[\textbf{20.}] (1) $f'(x)=\mathrm e^{x}-1$，故 $f$ 在 $(-\infty,0)$ 递减、在 $(0,+\infty)$ 递增；(2) $a>\dfrac1{\mathrm e}$
  \item[\textbf{21.}] (1) $\dfrac{x^{2}}9+y^{2}=1$；(2) 直线 $CD$ 过定点 $\left(\dfrac32,0\right)$
\end{itemize}

\gkpart{选做题}
\begin{itemize}[leftmargin=2.6em]
  \item[\textbf{22.}] (1) $C_1$ 是\textbf{以原点为圆心、半径 $1$ 的圆}；(2) 公共点为 $\left(\dfrac14,\dfrac14\right)$。
  \item[\textbf{23.}] (1) $f(x)=\begin{cases}-x-3,&x<-\dfrac13\\ 5x-1,&-\dfrac13\le x<1\\ x+3,&x\ge1\end{cases}$；(2) 解集为 $\left(-\infty,-\dfrac76\right)$。
\end{itemize}
